\chapter{Detailed computations}
\label{chap:detailed-computations}

This chapter provides complete computational details for the key examples of
Part~II (Examples), including all structure constants and explicit verifications.


\section{Heisenberg computations}
\label{sec:heisenberg-details}

\subsection{Bar complex: degrees 3 and 4}

\begin{computation}[Heisenberg bar complex: full degree 3]\label{comp:heisenberg-deg3-full}
The degree 3 component of $\B(\cH)$ is:
\[
\B_3(\cH) = \bigoplus_{m,n,p > 0} \bC \cdot [a_{-m}|a_{-n}|a_{-p}] \otimes \Omega^2(\overline{\Conf}_3)
\]

\emph{Basis for $\Omega^2(\overline{\Conf}_3)$.}
The three 2-forms $\eta_{12} \wedge \eta_{13}$, $\eta_{12} \wedge \eta_{23}$, and $\eta_{13} \wedge \eta_{23}$ (where $\eta_{ij} = d\log(z_i - z_j)$) satisfy the Arnold relation:
\[
\eta_{12} \wedge \eta_{23} + \eta_{23} \wedge \eta_{31} + \eta_{31} \wedge \eta_{12} = 0
\]
(with $\eta_{31} = \eta_{13}$, since $d\log(z_3 - z_1) = d\log(z_1 - z_3)$).  Therefore $\dim \Omega^2(\overline{\Conf}_3) = 2$, consistent with the Poincar\'{e} polynomial $P_t(C_3(\mathbb{C})) = 1 + 3t + 2t^2$.  A basis is $\{\eta_{12} \wedge \eta_{13},\; \eta_{13} \wedge \eta_{23}\}$, and the Arnold relation gives $\eta_{12} \wedge \eta_{23} = \eta_{12} \wedge \eta_{13} + \eta_{13} \wedge \eta_{23}$.

\emph{Differential computation.}
\[
d([a_{-m}|a_{-n}|a_{-p}] \otimes \eta_{12} \wedge \eta_{23}) = d_{\text{res}} + d_{\text{dR}}
\]

For $d_{\text{res}}$: We compute residues at $D_{12}$, $D_{23}$, $D_{13}$:
\begin{align*}
\Res_{D_{12}} &= \Res_{z_1 = z_2}\left[\frac{k}{(z_1-z_2)^2} \cdot \frac{d(z_1-z_2)}{z_1-z_2} \wedge \frac{d(z_2-z_3)}{z_2-z_3}\right] \\
&= \Res_{\epsilon \to 0}\left[\frac{k d\epsilon}{\epsilon^3} \wedge \eta_{23}\right] = 0
\end{align*}
(Triple pole gives zero residue.)

Similarly for $D_{23}$ and $D_{13}$: all residues vanish.

\emph{For $d_{\text{dR}}$.}
\[
d_{\text{dR}}(\eta_{12} \wedge \eta_{23}) = 0
\]
since $d(\eta_{ij}) = 0$ (logarithmic forms are closed on $\overline{\Conf}_3 \setminus \partial$).

Hence $d([a_{-m}|a_{-n}|a_{-p}] \otimes \eta_{12} \wedge \eta_{23}) = 0$.

All degree 3 elements with top form are cycles.
\end{computation}

\begin{computation}[Heisenberg: degree 4 differential]\label{comp:heisenberg-deg4}
For degree 4 with 4 tensor factors:
\[
\B_4(\cH) \ni [a_{-m}|a_{-n}|a_{-p}|a_{-q}] \otimes \omega_4
\]

The form space $\Omega^3(\overline{\Conf}_4)$ has dimension 6, generated by 
products of three $\eta_{ij}$.

The differential $d_{\text{res}}$ involves six collision divisors $D_{ij}$ 
($1 \leq i < j \leq 4$). At each divisor:
\[
\Res_{D_{ij}}[a_{-m_i} \cdot a_{-m_j} \otimes \omega] = \Res\left[\frac{k}{(z_i-z_j)^2} \otimes \omega\right]
\]

When $\omega$ contains exactly one factor of $\eta_{ij}$:
\[
\Res_{D_{ij}}\left[\frac{k}{(z_i-z_j)^2} \cdot \frac{d(z_i-z_j)}{z_i-z_j} \wedge \cdots\right] = 0
\]

When $\omega$ contains no factors of $\eta_{ij}$:
\[
\Res_{D_{ij}}\left[\frac{k}{(z_i-z_j)^2} \cdot \eta_{ik} \wedge \eta_{j\ell} \wedge \cdots\right]
\]
This gives a nontrivial contribution only if there is a simple pole, which requires 
the OPE to have a simple pole term (the Heisenberg OPE does not).

The entire degree 4 differential vanishes on elements with
maximal form degree, confirming the exterior algebra structure of homology.
\end{computation}


\subsection{Twisting morphism verification}

\begin{computation}[Maurer--Cartan through degree 4]\label{comp:heisenberg-mc-deg4}
We verify the Maurer--Cartan equation $d\tau + \tau \star \tau = 0$ in detail.

\emph{Setup.} The twisting morphism $\tau: \B(\cH) \to \cH$ is:
\[
\tau([a_{-n}]) = a_{-n}, \quad \tau(\text{higher}) = 0
\]

The convolution product $\tau \star \tau: \B(\cH) \to \cH$ is computed via:
\[
(\tau \star \tau)(x) = \mu(\tau \otimes \tau)\Delta(x)
\]
where $\Delta$ is the coproduct on $\B(\cH)$ and $\mu$ is the product on $\cH$.

\emph{On degree~2.}
\[
\Delta([a_{-m}|a_{-n}]) = 1 \otimes [a_{-m}|a_{-n}] + [a_{-m}] \otimes [a_{-n}] + [a_{-m}|a_{-n}] \otimes 1
\]
(the standard deconcatenation coproduct; the terms $1 \otimes (\cdots)$ and $(\cdots) \otimes 1$ give zero under $\tau \otimes \tau$ since $\tau$ vanishes on degree $\neq 1$).

Thus:
\[
(\tau \star \tau)([a_{-m}|a_{-n}]) = \tau([a_{-m}]) \cdot \tau([a_{-n}]) = a_{-m} \cdot a_{-n}
\]
(a single term, not symmetrized).

Meanwhile:
\[
d\tau([a_{-m}|a_{-n}]) = \tau(d[a_{-m}|a_{-n}]) = \tau(\text{residue terms}) = 0
\]
since residue terms involve the OPE coefficient $k$, not the generators themselves.

We have $d\tau = 0$ but $\tau \star \tau \neq 0$; the resolution is that the Heisenberg algebra is a commutative chiral algebra ($a_{(0)}b = 0$ for all generators), so the relevant bar complex is the \emph{Harrison subcomplex} (the commutative bar complex, where we impose shuffle antisymmetry).  In the Harrison complex:
\[
[a_{-m}|a_{-n}] + [a_{-n}|a_{-m}] = 0 \quad \text{(shuffle antisymmetry)}
\]
so $(\tau \star \tau)([a_{-m}|a_{-n}]) = a_{-m}a_{-n} - a_{-n}a_{-m} = [a_{-m}, a_{-n}]$,
which vanishes for $m + n \neq 0$.  When $m + n = 0$, the commutator
$[a_{-m}, a_m] = \kappa m \neq 0$; this non-vanishing is precisely the curvature
$m_0 = \kappa \cdot \mathbf{1}$ of the curved $A_\infty$ structure.

On the reduced bar complex, $d\tau = 0$ and
$\tau \star \tau = -m_0$ (the curved Maurer--Cartan equation), consistent with
the non-zero central extension.
\end{computation}


\section{Kac--Moody computations}
\label{sec:km-details}

\begin{computation}[$\widehat{\mathfrak{sl}}_3$ OPE]\label{comp:sl3-ope}
The Lie algebra $\mathfrak{sl}_3$ has basis:
\[
\{H_1, H_2, E_1, E_2, E_3, F_1, F_2, F_3\}
\]
where $E_3 = [E_1, E_2]$ and $F_3 = [F_2, F_1]$.

\emph{Cartan matrix.}
\[
A = \begin{pmatrix} 2 & -1 \\ -1 & 2 \end{pmatrix}
\]

\emph{Structure constants.}
\begin{align*}
[H_i, E_j] &= A_{ij} E_j, \quad [H_i, F_j] = -A_{ij} F_j \\
[E_i, F_j] &= \delta_{ij} H_i \quad (i,j \in \{1,2\}) \\
[E_1, E_2] &= E_3, \quad [F_2, F_1] = F_3 \\
[E_1, E_3] &= 0 = [E_2, E_3], \quad \text{similarly for } F
\end{align*}

The normalized trace form is $(H_i, H_j) = \mathrm{tr}_{\mathrm{fund}}(H_i H_j) = A_{ij}$, $(E_i, F_j) = \delta_{ij}$.
(The Killing form is $\kappa = 2h^\vee \cdot \mathrm{tr}_{\mathrm{fund}} = 6\,\mathrm{tr}_{\mathrm{fund}}$ for $\mathfrak{sl}_3$;
the OPEs below use the normalized form with level $k$ measured in units of the long root.)

The affine OPEs at level $k$ are:
\begin{align*}
H_i(z) H_j(w) &\sim \frac{kA_{ij}}{(z-w)^2} \\
H_i(z) E_j(w) &\sim \frac{A_{ij} E_j(w)}{z-w} \\
E_i(z) F_j(w) &\sim \frac{k\delta_{ij}}{(z-w)^2} + \frac{\delta_{ij} H_i(w)}{z-w} \\
E_1(z) E_2(w) &\sim \frac{E_3(w)}{z-w} \\
E_1(z) E_3(w) &\sim 0 \quad \text{(since $[E_1, E_3] = 0$ in $\mathfrak{sl}_3$)}
\end{align*}
\end{computation}

\begin{computation}[$\widehat{\mathfrak{sl}}_3$ bar differential]\label{comp:sl3-bar}
The degree-2 differential is as follows.

With constant form:
\begin{align*}
d[H_{1,-m}|E_{1,-n}] &= 2[E_{1,-(m+n)}] \\
d[H_{1,-m}|E_{2,-n}] &= -[E_{2,-(m+n)}] \\
d[E_{1,-m}|E_{2,-n}] &= [E_{3,-(m+n)}] \\
d[E_{1,-m}|F_{1,-n}] &= [H_{1,-(m+n)}] + km\delta_{m+n,0}[1]
\end{align*}

With $\eta_{12}$: All differentials vanish (no simple poles when paired with $d\log$).

The degree-3 differential encodes the Serre relations.

The bar complex encodes the Serre relations:
\[
[E_{1,-\ell}|E_{1,-m}|E_{2,-n}] - 2[E_{1,-\ell}|E_{2,-m}|E_{1,-n}] + [E_{2,-\ell}|E_{1,-m}|E_{1,-n}]
\]

The differential of this combination gives the Serre identity:
\[
d(\text{above}) = E_1 E_1 E_2 - 2 E_1 E_2 E_1 + E_2 E_1 E_1 = (\text{ad}\, E_1)^2(E_2) = 0 \quad \text{in } U(\hat{\mathfrak{sl}}_3)
\]
\end{computation}


\subsection{Acyclicity verification}

\begin{computation}[Kac--Moody acyclicity at generic level]\label{comp:km-acyclic}
We verify acyclicity of $\B(\hat{\mathfrak{sl}}_2) \otimes_\tau V_k$ for generic $k$.

\emph{Filtration.} Define $F_p$ by the PBW degree (number of generators applied to vacuum).

\emph{Associated graded.}
\[
\mathrm{gr}_F(\B(\hat{\mathfrak{sl}}_2) \otimes_\tau V_k) \cong \B(\mathfrak{sl}_2[t^{-1}]) \otimes S(\mathfrak{sl}_2[t^{-1}])
\]
the bar complex of the loop algebra tensored with the symmetric algebra.

\emph{$E_1$ page.} The homology is:
\[
E_1^{p,q} = H_p(\mathfrak{sl}_2; S^q(\mathfrak{sl}_2[t^{-1}]))
\]
Lie algebra homology with polynomial coefficients.

\emph{Vanishing.} For $\mathfrak{sl}_2$:
\[
H_n(\mathfrak{sl}_2; M) = 0 \quad \text{for } n > 0 \text{ and } M \text{ a rational } \mathfrak{sl}_2\text{-module}
\]
by the Whitehead lemmas (semisimplicity).

\emph{Convergence.} The spectral sequence collapses at $E_1$, giving:
\[
H_n(\B(\hat{\mathfrak{sl}}_2) \otimes_\tau V_k) = \begin{cases}
\bC & n = 0 \\
0 & n > 0
\end{cases}
\]
confirming acyclicity.

\emph{At special levels.} When $k = -h^\vee = -2$ (critical level for $\mathfrak{sl}_2$), 
the vacuum module becomes reducible, and new homology classes appear, corresponding 
to the center of $U(\hat{\mathfrak{sl}}_2)_{-2}$.
\end{computation}


\section{W-algebra computations}
\label{sec:w-algebra-details}

\subsection{Virasoro bar complex through degree 3}

\begin{computation}[Virasoro OPE products]\label{comp:virasoro-ope}
\index{Virasoro algebra!bar complex}
The Virasoro algebra has a single strong generator $T$ (conformal weight~$2$) with OPE:
\begin{align*}
T(z)\,T(w) &= \frac{c/2}{(z-w)^4} + \frac{2T(w)}{(z-w)^2} + \frac{\partial T(w)}{z-w} + \normord{TT}(w) + O(z-w)
\end{align*}
The $n$-th products (singular terms):
\begin{center}
\begin{tabular}{clll}
$n$ & $T_{(n)}T$ & Pole order & Role \\
\hline
$3$ & $c/2$ & $(z\!-\!w)^{-4}$ & curvature \\
$2$ & $0$ & $(z\!-\!w)^{-3}$ & (absent) \\
$1$ & $2T$ & $(z\!-\!w)^{-2}$ & conformal weight \\
$0$ & $\partial T$ & $(z\!-\!w)^{-1}$ & translation
\end{tabular}
\end{center}

\emph{Key difference from Kac--Moody:} the affine Lie algebra $\widehat{\mathfrak{g}}_k$ has OPE poles up to order~$2$, while the Virasoro has a quartic pole.  This extra pole produces the curvature element $m_0 = (c/2)\cdot\mathbf{1}$ in the $A_\infty$ structure (cf.\ Computation~\ref{comp:virasoro-bar-diff} below).
\end{computation}

\begin{computation}[Virasoro vacuum module: low-weight states]\label{comp:virasoro-vacuum}
The augmentation ideal $\bar{V} = V/\mathbb{C}\cdot|0\rangle$ of the Virasoro vacuum module has a basis of states created by $L_{-n}$ (for $n \geq 2$) applied to the vacuum:
\begin{center}
\begin{tabular}{c|cccccc}
weight $h$ & 2 & 3 & 4 & 5 & 6 & $h$ \\
\hline
$\dim \bar{V}_h$ & 1 & 1 & 2 & 2 & 4 & $p_{\geq 2}(h)$ \\
basis & $T$ & $\partial T$ & $\partial^2 T,\, \normord{TT}$ & $2$ & $4$ &
\end{tabular}
\end{center}
where $p_{\geq 2}(h) = p(h) - p(h-1)$ denotes the number of partitions of~$h$ into parts $\geq 2$, with generating function $\prod_{n=2}^{\infty}(1-q^n)^{-1}$.  The constraint arises from $L_{-1}|0\rangle = 0$ (translation invariance).
\end{computation}

\begin{computation}[Bar differential: degree~2 $\to$ degree~1]\label{comp:virasoro-bar-diff}
\index{bar differential!Virasoro}
The bar complex $\bar{B}(\mathrm{Vir}_c)$ has differential $D$ whose action on degree-2 elements extracts the full OPE data (all singular terms), as in the Kac--Moody formula~\eqref{eq:bar-diff-general}:
\begin{align}
D(T \otimes T \otimes \eta_{12}) &= T_{(3)}T \cdot |0\rangle + T_{(2)}T + T_{(1)}T + T_{(0)}T \notag \\
&= \frac{c}{2}\cdot|0\rangle + 0 + 2T + \partial T \label{eq:vir-bar-diff-TT}
\end{align}
The three contributions correspond to:
\begin{itemize}
\item $c/2\cdot|0\rangle \in \bar{B}^0$: the \emph{vacuum term} from the quartic pole (curvature)
\item $2T \in \bar{B}^1$: the \emph{conformal weight term} from the double pole
\item $\partial T \in \bar{B}^1$: the \emph{translation term} from the simple pole
\end{itemize}

\emph{Comparison with Kac--Moody:}
\begin{center}
\begin{tabular}{lll}
\textbf{Algebra} & \textbf{$D(v \otimes w \otimes \eta)$} & \textbf{Highest pole} \\
\hline
Heisenberg $\mathcal{H}_\kappa$ & $\kappa\cdot|0\rangle$ & double ($z^{-2}$) \\
$\widehat{\mathfrak{sl}}_{2,k}$ & $k(J^a,J^b)|0\rangle + f^{ab}{}_c J^c$ & double ($z^{-2}$) \\
Virasoro $\mathrm{Vir}_c$ & $(c/2)|0\rangle + 2T + \partial T$ & quartic ($z^{-4}$)
\end{tabular}
\end{center}

At weight $h = 5$, the bar differential on mixed elements $T \otimes \partial T$ and $\partial T \otimes T$:
\begin{align}
D(T \otimes \partial T \otimes \eta_{12}) &= T_{(3)}(\partial T) \cdot |0\rangle + T_{(1)}(\partial T) + T_{(0)}(\partial T) \notag \\
&= 0 + 3\partial T + \partial^2 T \label{eq:vir-bar-TdT}
\end{align}
using $T_{(3)}(\partial T) = \partial(T_{(3)}T) + 3 T_{(2)}T = 0 + 0 = 0$, $T_{(1)}(\partial T) = \partial(T_{(1)}T) + T_{(0)}T = 2\partial T + \partial T = 3\partial T$, and $T_{(0)}(\partial T) = \partial(T_{(0)}T) = \partial^2 T$.
\end{computation}

\begin{computation}[Curvature element and $c=0$ special value]\label{comp:virasoro-curvature}
The curvature element of the Virasoro bar complex is:
\begin{equation}\label{eq:vir-curvature}
m_0 = \frac{c}{2} \in \bar{B}^0 = \mathbb{C}
\end{equation}
arising from the quartic pole $T_{(3)}T = c/2$.

The $A_\infty$ relation $m_1^2 = [m_0, -]$ gives:
\begin{equation}
m_1^2(v) = m_2(m_0, v) - m_2(v, m_0)
\end{equation}
for any $v \in \bar{V}$.  The curvature controls the failure of $m_1$ (the linearized bar differential) to be a differential: $m_1^2 \neq 0$ unless $c = 0$.

\emph{Three regimes.}
\begin{itemize}
\item $c = 0$: Uncurved.  The bar complex is a genuine cochain complex ($D^2 = 0$ on $\bar{V}$).  The Koszul dual is $\mathrm{Vir}_0^! \simeq \mathrm{Vir}_0$ (Theorem~\ref{thm:virasoro-self-duality}).  The $TT$ OPE is purely quadratic ($2T/(z-w)^2 + \partial T/(z-w)$), and the Koszul duality is of classical type.
\item $c \neq 0$, generic: Curved $A_\infty$.  The completed bar construction is needed (Appendix~\ref{app:nilpotent-completion}).  The Koszul dual is $\mathrm{Vir}_c^! \simeq \mathrm{Vir}_{26-c}$ (Example~\ref{ex:virasoro-koszul-dual}), where the shift $c \mapsto 26-c$ arises from the Serre duality pairing on $\overline{C}_{n+1}(X)$ interacting with the quartic pole data.
\item $c = 26$: The dual Virasoro $\mathrm{Vir}_{26}^! \simeq \mathrm{Vir}_0$ is the uncurved algebra.  Equivalently, the bosonic string BRST complex at $c=26$ is governed by the bar complex of $\mathrm{Vir}_{26}$ (Theorem~\ref{thm:virasoro-string}).
\end{itemize}

\emph{The curvature--central charge dictionary:} the curvature $m_0 = c/2$ is an invariant of the chiral algebra that determines the qualitative structure of the bar complex.  For the Virasoro, it is the \emph{only} such invariant (since $T$ is the sole generator), so the full bar complex structure is controlled by $c$ alone.  This is the simplest instance of the general principle that curvature = quantum correction (Theorem~\ref{thm:genus-universality}).
\end{computation}

\begin{computation}[$\mathcal{W}_3$ OPE coefficients]\label{comp:w3-coeff}
The $\mathcal{W}_3$ algebra has generators $T$ (spin 2) and $W$ (spin 3).

\emph{Full $W(z)W(w)$ OPE.}
\begin{align*}
W(z)W(w) &= \frac{c/3}{(z-w)^6} + \frac{2T(w)}{(z-w)^4} + \frac{\partial T(w)}{(z-w)^3} \\
&\quad + \frac{\frac{3}{10}\partial^2 T + \frac{16}{22+5c}\Lambda}{(z-w)^2} \\
&\quad + \frac{\frac{1}{15}\partial^3 T + \frac{8}{22+5c}\partial\Lambda}{z-w} + \text{regular}
\end{align*}
where:
\[
\Lambda = \normord{TT} - \frac{3}{10}\partial^2 T
\]

\emph{Verification of Jacobi.} The $T$-$W$-$W$ Jacobi identity gives:
\[
[L_m, [W_n, W_p]] + \text{cyclic} = 0
\]
Using $[L_m, W_n] = (2m - n)W_{m+n}$:
\begin{align*}
[L_m, [W_n, W_p]] &= [L_m, \text{(sum of } L_{n+p}, \Lambda_{n+p}, \text{etc.)}] \\
&= (m - n - p)[\cdots] + \cdots
\end{align*}
The coefficient $\frac{16}{22+5c}$ is uniquely determined by requiring Jacobi to hold.

\emph{Central charge values.}
\begin{itemize}
\item $c = -2$: A non-unitary $W_3$ algebra.  The DS parametrization gives $k = -3/2$, corresponding to $(p,q) = (3,2)$ with $\gcd(3,2) = 1$, but $W_3$ minimal models require $p, q \geq 3$, so $c = -2$ lies outside the minimal model series.  (Not to be confused with symplectic fermions, which also have $c=-2$ but are a Virasoro, not $W_3$, vertex algebra.)
\item $c = 0$: Trivial/degenerate $W_3(4,3)$ model (not to be confused with symplectic fermions, which have $c=-2$).
\item $c = 2$: Free field (Heisenberg + bc ghosts).
\item $c \to \infty$: Classical limit ($\mathcal{W}_3^{\mathrm{cl}}$).
\end{itemize}
\end{computation}


\begin{computation}[DS reduction for $\mathcal{W}_3$]\label{comp:ds-w3}
Starting from $\hat{\mathfrak{sl}}_3$ at level $k$:

\emph{BRST complex.}
\[
C^\bullet = V_k(\mathfrak{sl}_3) \otimes \mathrm{Cl}(\mathfrak{n}((t)) \oplus \mathfrak{n}^*((t)))
\]
where $\mathfrak{n} = \bC E_1 \oplus \bC E_2 \oplus \bC E_3$ is the nilpotent radical ($E_3 = [E_1, E_2]$) and $\mathrm{Cl}$ denotes the Clifford algebra (ghost--anti-ghost system).

\emph{Ghost fields.} $c_1, c_2, c_3$ (fermionic) with $|c_i| = 1$.

\emph{BRST differential.}
\begin{align*}
Q &= \sum_n (E_{1,n} - \chi_1\delta_{n,0})c_{1,-n} + \sum_n (E_{2,n} - \chi_2\delta_{n,0})c_{2,-n} \\
&\quad + \sum_n E_{3,n} c_{3,-n} + \sum_{m,n} c_{1,m}c_{2,n}b_{3,-(m+n)} + \cdots
\end{align*}
where $\chi_1, \chi_2$ determine the nilpotent element $f = \chi_1 F_1 + \chi_2 F_2$.

\emph{Principal nilpotent.} $\chi_1 = \chi_2 = 1$.

\emph{$H^0(Q)$.} Generated by:
\begin{itemize}
\item $T = $ Sugawara tensor (survives reduction)
\item $W = $ new spin-3 generator from reduction
\end{itemize}

\emph{Central charge formula.}
\[
c_{\mathcal{W}_3}(k) = 2 - \frac{24(k+2)^2}{k+3}
\]
\end{computation}


\subsection{\texorpdfstring{$\mathcal{W}_3$ bar complex: degree-2 computation}{W-3 bar complex: degree-2 computation}}

\begin{computation}[Complete $n$-th products for $\mathcal{W}_3$]\label{comp:w3-nthproducts}
\index{W3 algebra@$\mathcal{W}_3$!$n$-th products}
The $\mathcal{W}_3$ algebra has generators $T$ (weight~$2$) and $W$ (weight~$3$).
We list all singular $n$-th products ($a_{(n)}b$ for $n \geq 0$ with nonzero coefficient)
for each ordered generator pair.

\emph{$T \times T$ (quartic pole).}
\begin{align*}
T_{(3)}T &= c/2, & T_{(2)}T &= 0, &
T_{(1)}T &= 2T, & T_{(0)}T &= \partial T.
\end{align*}

\emph{$T \times W$ (double pole).}
Since $W$ is primary of weight~$3$, the $TW$ OPE has at most a double pole:
\begin{align*}
T_{(1)}W &= 3W, & T_{(0)}W &= \partial W.
\end{align*}
All $T_{(n)}W = 0$ for $n \geq 2$.

\emph{$W \times T$ (double pole).}
By the vertex algebra skew-symmetry formula
$b_{(n)}a = \sum_{j \geq 0} \frac{(-1)^{n+j+1}}{j!}\,\partial^j(a_{(n+j)}b)$:
\begin{align*}
W_{(1)}T &= T_{(1)}W = 3W, \\
W_{(0)}T &= -T_{(0)}W + \partial(T_{(1)}W) = -\partial W + 3\partial W = 2\partial W.
\end{align*}
The asymmetry $W_{(0)}T = 2\partial W \neq \partial W = T_{(0)}W$ is a consequence
of the derivative correction in the skew-symmetry formula.  All $W_{(n)}T = 0$ for $n \geq 2$.

\emph{$W \times W$ (sixth-order pole).}
\begin{center}
\begin{tabular}{clll}
$n$ & $W_{(n)}W$ & Pole order & Weight \\
\hline
$5$ & $c/3$ & $(z\!-\!w)^{-6}$ & $0$ \\
$4$ & $0$ & $(z\!-\!w)^{-5}$ & -- \\
$3$ & $2T$ & $(z\!-\!w)^{-4}$ & $2$ \\
$2$ & $\partial T$ & $(z\!-\!w)^{-3}$ & $3$ \\
$1$ & $\tfrac{3}{10}\partial^2 T + \tfrac{16}{22+5c}\Lambda$ & $(z\!-\!w)^{-2}$ & $4$ \\
$0$ & $\tfrac{1}{15}\partial^3 T + \tfrac{8}{22+5c}\partial\Lambda$ & $(z\!-\!w)^{-1}$ & $5$
\end{tabular}
\end{center}
where $\Lambda = \normord{TT} - \frac{3}{10}\partial^2 T$ (Definition~\ref{def:lambda-complete}).
The vanishing $W_{(4)}W = 0$ occurs because $\dim \bar{V}_1 = 0$ in the vacuum module
($L_{-1}|0\rangle = 0$), so there is no weight-$1$ state to appear at the fifth-order pole.
\end{computation}


\begin{computation}[$\mathcal{W}_3$ bar differential on degree-2 elements]\label{comp:w3-bar-degree2}
\index{bar differential!$\mathcal{W}_3$}
The bar differential $D \colon \bar{B}^2(\mathcal{W}_3) \to \bar{B}^0 \oplus \bar{B}^1$
extracts all singular OPE data from each generator pair.
We compute $D(a \otimes b \otimes \eta_{12})$ for all four cases.

\emph{(i)} $D(T \otimes T \otimes \eta_{12})$
(from Computation~\ref{comp:virasoro-bar-diff}):
\[
D(T \otimes T \otimes \eta_{12}) = \frac{c}{2}\cdot|0\rangle + 2T + \partial T
\]

\emph{(ii)} $D(T \otimes W \otimes \eta_{12})$:
\[
D(T \otimes W \otimes \eta_{12}) = T_{(1)}W + T_{(0)}W = 3W + \partial W
\]
No vacuum contribution (the $TW$ OPE has no pole beyond order~$2$).

\emph{(iii)} $D(W \otimes T \otimes \eta_{12})$:
\[
D(W \otimes T \otimes \eta_{12}) = W_{(1)}T + W_{(0)}T = 3W + 2\partial W
\]
The asymmetry with~(ii) comes from skew-symmetry
(Computation~\ref{comp:w3-nthproducts}).

\emph{(iv)} $D(W \otimes W \otimes \eta_{12})$:
\begin{align*}
D(W \otimes W \otimes \eta_{12})
&= W_{(5)}W \cdot |0\rangle + W_{(3)}W + W_{(2)}W + W_{(1)}W + W_{(0)}W \\
&= \frac{c}{3}\cdot|0\rangle + 2T + \partial T
  + \Bigl[\frac{3}{10}\partial^2 T + \frac{16}{22+5c}\Lambda\Bigr] \\
&\quad + \Bigl[\frac{1}{15}\partial^3 T + \frac{8}{22+5c}\partial\Lambda\Bigr]
\end{align*}
(The $n=4$ term vanishes: $W_{(4)}W = 0$.)
This is the most complex bar differential in the manuscript: it produces components
in every weight from~$0$ (vacuum) through~$5$ ($\partial\Lambda$).

\emph{Summary of vacuum contributions:}
\begin{center}
\begin{tabular}{lcc}
\textbf{Generator pair} & \textbf{Vacuum term} & \textbf{Max weight in $\bar{B}^1$} \\
\hline
$T \otimes T$ & $c/2$ & $3$ ($\partial T$) \\
$T \otimes W$ & $0$ & $4$ ($\partial W$) \\
$W \otimes T$ & $0$ & $4$ ($2\partial W$) \\
$W \otimes W$ & $c/3$ & $5$ ($\partial\Lambda$)
\end{tabular}
\end{center}
Only the ``diagonal'' pairs $T \otimes T$ and $W \otimes W$ produce vacuum leakage.
\end{computation}


\begin{computation}[$\mathcal{W}_3$ curvature and dual level verification]\label{comp:w3-curvature-dual}
\index{curvature!$\mathcal{W}_3$ bar complex}
The $\mathcal{W}_3$ bar complex has vacuum leakage from two independent channels:
\begin{enumerate}
\item \emph{Virasoro sector}: $T_{(3)}T = c/2$ (quartic pole), contributing $m_0^{(T)} = c/2$.
\item \emph{$W$-sector}: $W_{(5)}W = c/3$ (sixth-order pole), contributing $m_0^{(W)} = c/3$.
\end{enumerate}
The ratio $m_0^{(W)}/m_0^{(T)} = 2/3$ is independent of the level and equals the
ratio of central term normalizations in the $WW$ and $TT$ OPEs
($c/3$ vs.\ $c/2$).

Both leakages vanish simultaneously if and only if $c = 0$.
The full $\mathcal{W}_3$ curvature is proportional to $(k + h^\vee) = (k+3)$
(Theorem~\ref{thm:w-bar-curvature}), which vanishes at the critical level
$k = -3$.  These are \emph{distinct} conditions:
\begin{itemize}
\item $c = 0$ occurs at rational levels $k = -5/3$ and $k = -9/4$
  (from solving $12(k+2)^2 = k+3$, i.e., $12u^2 - u - 1 = 0$ where $u = k+2$).
\item $k = -3$ is the critical level where the Sugawara tensor is undefined and
  the DS construction degenerates.
\end{itemize}
Similarly, for the Virasoro ($\mathfrak{sl}_2$), $c = 0$ occurs at
$k = -1/2$ and $k = -4/3$ (both rational), while critical level is $k = -2$.
See Theorem~\ref{thm:w-bar-curvature} for the full discussion.

At the dual level $k' = -k-6$, the dual central charge is $c' = c(k') = 2 - 24(k'+2)^2/(k'+3) = 2 + 24(k+4)^2/(k+3)$.
The sum:
\[
c + c' = 4 + 24\cdot\frac{(k+4)^2 - (k+2)^2}{k+3}
       = 4 + 24\cdot\frac{4(k+3)}{k+3} = 100
\]
confirming the general complementarity $c(\mathcal{W}^k) + c(\mathcal{W}^{k'})
= 2r + 4h^\vee d = 2(2) + 4(3)(8) = 100$
(Theorem~\ref{thm:central-charge-complementarity}).

At level $k'$, the vacuum leakages are $m_0^{(T)} = c'/2$ and $m_0^{(W)} = c'/3$.
Since $c' = 100 - c$, the sum of curvatures satisfies:
\[
m_0^{(T)}(k) + m_0^{(T)}(k') = \frac{c}{2} + \frac{100-c}{2} = 50
\]
The curvature of one algebra determines the curvature of its dual.  When
$c = 0$ (uncurved at $k = -5/3$ or $-9/4$), the dual has $c' = 100$
(strongly curved).  The involution $k \mapsto -k-6$ has fixed point $k = -3$
(the critical level), where $c$ diverges and the DS construction degenerates;
the complementarity formula $c + c' = 100$ governs the curvature trade-off
for all non-critical levels but has no $\mathcal{W}$-algebraic meaning at
the critical level itself.

The composite field coefficient at level $k$ is $\alpha(c) = 16/(22+5c)$, and
at dual level $\alpha(c') = 16/(22+5(100-c)) = 16/(522-5c)$.
These differ (the dual algebra has a different OPE structure), but both are
determined by a single parameter $c$.  In the large-level limit $k \to \infty$
(where $c \to -\infty$), both $\alpha(c) \to 0^-$ and
$\alpha(c') \to 0^+$: both algebras become classical simultaneously.
\end{computation}


\section{Yangian bar complex details}
\label{sec:yangian-details}

\begin{computation}[Yangian $Y(\mathfrak{sl}_2)$ relations]\label{comp:yangian-sl2}
Generators: $e^{(r)}, f^{(r)}, h^{(r)}$ for $r \geq 0$.

\emph{Level~0} (= $\mathfrak{sl}_2$):
\[
[h^{(0)}, e^{(0)}] = 2e^{(0)}, \quad [h^{(0)}, f^{(0)}] = -2f^{(0)}, \quad [e^{(0)}, f^{(0)}] = h^{(0)}
\]

\emph{Level~1.}
\begin{align*}
[h^{(0)}, e^{(1)}] &= 2e^{(1)} \\
[h^{(1)}, e^{(0)}] &= 2e^{(1)} \\
[e^{(0)}, f^{(1)}] &= h^{(1)} \\
[e^{(1)}, f^{(0)}] &= h^{(1)} \\
[e^{(1)}, f^{(1)}] - [e^{(0)}, f^{(2)}] &= \frac{1}{4}\{h^{(0)}, h^{(1)}\}
\end{align*}

\emph{Yangian current.}
\[
e(u) = \sum_{r \geq 0} e^{(r)} u^{-r-1}
\]

\emph{RTT presentation.}
\[
R(u-v) T_1(u) T_2(v) = T_2(v) T_1(u) R(u-v)
\]
with $R(u) = 1 - \frac{\hbar P}{u}$ and $T(u) = \begin{pmatrix} k_+(u) & e(u) \\ f(u) & k_-(u) \end{pmatrix}$ (matching the sign convention of Definition~\ref{def:yangian-rtt}).
\end{computation}

\begin{computation}[Yangian bar complex]\label{comp:yangian-bar-full}
\emph{Degree~1 generators.} $[e^{(r)}_{-n}], [f^{(r)}_{-n}], [h^{(r)}_{-n}]$ for $r \geq 0$, $n > 0$.

\emph{Degree~2 differential with R-matrix.}

The Yangian $R$-matrix $R(u) = 1 - \hbar P/u$ (where $P$ is the permutation and $u$ is the spectral parameter; cf.\ Definition~\ref{def:yangian-rtt}) deforms the bar differential.  At the level of the Yangian generating series $T^a(u) = \sum_{r \geq 0} T^{a,(r)} u^{-r-1}$, the deformed differential on degree-2 elements is:
\[
d^R[a^{(r)}|b^{(s)}] = \sum_{t=0}^{r+s} [a^{(r)}, b^{(s)}]^{(t)} + \text{(R-matrix corrections at level } r+s+1 \text{)}
\]
where the Yangian brackets $[a^{(r)}, b^{(s)}]^{(t)}$ encode the RTT relations.

The \emph{loop modes} $a^{(r)}_{-n}$ (combining Yangian level $r$ with affine mode $-n$) carry two independent gradings.  The bar differential must respect both: the Yangian level $r$ governs the $R$-matrix corrections, while the affine mode $n$ governs the Lie bracket.  For the simplest case ($r = s = 0$):
$$d^R[e^{(0)}_{-m}|f^{(0)}_{-n}] = [h^{(0)}_{-(m+n)}] + O(\text{level } \geq 1)$$
where the leading term is the standard affine Lie algebra bracket and higher-level corrections involve $T^{(r)}$ for $r \geq 1$.

\emph{Cohomology.} Related to Yangian cohomology $H^*(Y(\mathfrak{sl}_2))$.
\end{computation}


\subsection{\texorpdfstring{$q$-deformed computations}{q-deformed computations}}\label{sec:q-deformed-details}

\begin{computation}[$q$-Heisenberg commutator expansion]\label{comp:q-comm}
The $q$-commutator:
\[
[a_m, a_n]_q = q^{\mathrm{sgn}(n-m)/2} a_m a_n - q^{\mathrm{sgn}(m-n)/2} a_n a_m
\]

\emph{For $m < n$.}
\[
[a_m, a_n]_q = q^{1/2} a_m a_n - q^{-1/2} a_n a_m
\]

\emph{For $m > n$.}
\[
[a_m, a_n]_q = q^{-1/2} a_m a_n - q^{1/2} a_n a_m = -[a_n, a_m]_q
\]
(antisymmetry preserved).

\emph{For $m = -n$.}
\[
[a_m, a_{-m}]_q = [m]_q
\]
where $[m]_q = \frac{q^m - q^{-m}}{q - q^{-1}}$.

\emph{Expansion as $q \to 1$.}
\[
[a_m, a_n]_q = [a_m, a_n] + \frac{\ln q}{2}\mathrm{sgn}(n-m)\{a_m, a_n\} + O((\ln q)^2)
\]
The first-order correction involves the anticommutator.
\end{computation}


\subsection{Roots of unity phenomena}

\begin{computation}[$q$-Heisenberg at roots of unity]\label{comp:q-root-unity}
Let $q = e^{2\pi i/N}$ be a primitive $N$-th root of unity.

\emph{Central elements.} The element $a_0^N$ becomes central:
\[
[a_0^N, a_m] = 0 \quad \text{for all } m
\]
(since $q^{mN} = 1$).

\emph{New relations.}
\[
a_m^N = 0 \quad \text{for } m \neq 0 \text{ (in certain quotients)}
\]

\emph{Bar complex modification.} New generators:
\[
[a_0^N], [a_1^N], \ldots
\]
with modified differential reflecting the truncation.

\emph{New homology.} Classes corresponding to the restricted Lie algebra structure.
\end{computation}


\section{Higher genus formulas}
\label{sec:higher-genus-details}

\subsection{Genus 1 (torus) formulas}

\begin{computation}[Heisenberg on torus]\label{comp:heisenberg-torus}
On the torus $E_\tau = \bC/(\bZ + \tau\bZ)$:

\emph{Green's function.}
\[
G(z, w | \tau) = -\log|\theta_1(z-w|\tau)| + \frac{\pi(\mathrm{Im}(z-w))^2}{\mathrm{Im}(\tau)}
\]
where $\theta_1$ is the odd Jacobi theta function.

\emph{OPE on torus.}
\[
J(z) J(w) = \kappa\,\wp(z-w|\tau) + \text{regular}
\]
where $\wp$ is the Weierstrass $\wp$-function with double pole at $z = w$.

\emph{Mode expansion.}
\[
J(z) = \sum_{n \in \bZ} a_n e^{2\pi i n z}
\]

\emph{Commutation relations.}
\[
[a_m, a_n] = \kappa\, m\,\delta_{m+n,0}
\]
The commutation relations are independent of $\tau$; the genus-1 data enters through the propagator $\wp(z|\tau) - \frac{\pi^2 E_2(\tau)}{3}$, not through the algebra structure.

\emph{Chiral partition function} (character of the Fock space):
\[
Z_{\mathrm{chiral}}(\tau) = \mathrm{tr}_{\mathcal{F}}(q^{L_0 - 1/24}) = \frac{1}{\eta(\tau)} = q^{-1/24} \prod_{n=1}^\infty \frac{1}{1 - q^n}
\]
where $q = e^{2\pi i \tau}$.  The full (non-chiral) partition function is $Z(\tau) = |\eta(\tau)|^{-2}$.
\end{computation}


\begin{computation}[Bar complex on higher genus]\label{comp:bar-genus-g}
For genus $g$ curve $\Sigma_g$:

\emph{Configuration space.} $\Conf_n(\Sigma_g)$ with compactification
$\overline{\Conf}_n(\Sigma_g)$.

\emph{Cohomology.}
$H^*(\Conf_n(\Sigma_g))$ is computed by the Cohen--Taylor (or Totaro) spectral sequence, converging from $E_1 = H^*(\Sigma_g^n) \otimes H^*(\text{OS}_n)$ (the Orlik--Solomon algebra of logarithmic forms $\omega_{ij}$ modulo Arnold relations) together with the mixed relations $\omega_{ij} \cdot [\text{diagonal}]_{ij} = 0$.  In particular, $H^1(\Sigma_g) \cong \mathbb{C}^{2g}$ contributes new generators.

\emph{New bar complex generators.} For each handle ($A_i$, $B_i$ cycle):
\[
[p_i], [q_i] \quad (i = 1, \ldots, g)
\]
representing the zero-modes from the period integrals.

\emph{Modified differential.}
\[
d[p_i | q_j] = k\delta_{ij} \cdot [\text{fundamental class}]
\]
reflecting the symplectic pairing $\int_{A_i} \omega \cdot \int_{B_j} \omega = \delta_{ij}$.

\emph{Homology.}
\[
H_n(\B(\cH_g)) = \wedge^n(V^* \oplus \bC^{2g})
\]
with the additional $2g$ generators from the handles.
\end{computation}


\section{Structure theorems from bar complex computations}
\label{sec:bar-structure-theorems}
\index{bar complex!structure theorems}

We collect the principal structural results established by the preceding
computations, stated as formal propositions.

\begin{proposition}[Heisenberg bar complex: maximal-form cycles; \ClaimStatusProvedHere]
\label{prop:heisenberg-maximal-form-cycles}
\index{Heisenberg algebra!maximal-form cycles}
For the Heisenberg algebra $\cH_\kappa$ at any level~$\kappa$, every element
of $\B_n(\cH_\kappa)$ with maximal form degree (i.e., tensored with a form
$\omega \in \Omega^{n-1}(\overline{\Conf}_n)$) is a cycle of the bar
differential, for all $n \geq 2$.
\end{proposition}

\begin{proof}
The bar differential has two components: the residue differential $d_{\mathrm{res}}$
(from OPE singularities at collision divisors) and the de~Rham differential
$d_{\mathrm{dR}}$ (from exterior differentiation on configuration spaces).

\emph{The de~Rham component.}
For maximal-degree forms $\omega \in \Omega^{n-1}(\overline{\Conf}_n)$,
$d_{\mathrm{dR}}(\omega) = 0$ since logarithmic forms $\eta_{ij} = d\log(z_i - z_j)$
are closed on $\overline{\Conf}_n \setminus \partial$.

\emph{The residue component.}
The Heisenberg OPE $a(z)\,a(w) = \kappa/(z-w)^2$ has a double pole
and no simple pole: $a_{(1)}a = \kappa$ (scalar) and $a_{(0)}a = 0$.
At each collision divisor~$D_{ij}$, we must evaluate
$\Res_{D_{ij}}(\kappa/(z_i - z_j)^2 \cdot \omega)$.
Two cases arise:
\begin{enumerate}
\item If $\omega$ contains the factor $\eta_{ij} = d(z_i-z_j)/(z_i-z_j)$,
  then the integrand is $\kappa\,d(z_i-z_j)/(z_i-z_j)^3 \wedge (\text{other factors})$.
  This triple pole has zero residue (the residue extracts the coefficient
  of $(z_i-z_j)^{-1}\,d(z_i-z_j)$, which is absent).
\item If $\omega$ does not contain $\eta_{ij}$, then the integrand
  $\kappa/(z_i-z_j)^2 \cdot \omega$ has a double pole with no
  $d(z_i-z_j)/(z_i-z_j)$ factor; the contour integral around $D_{ij}$
  vanishes.
\end{enumerate}
In both cases the residue is zero, so $d_{\mathrm{res}} = 0$ on
maximal-form elements.
Computations~\ref{comp:heisenberg-deg3-full} and~\ref{comp:heisenberg-deg4}
verify this explicitly at $n = 3$ and $n = 4$.
\end{proof}

\begin{remark}\label{rem:maximal-form-mechanism}
The vanishing mechanism is specific to the Heisenberg: the OPE has
only a double pole and no simple pole.  For the Virasoro
(quartic pole $T_{(3)}T = c/2$), the maximal-form bar differential
does \emph{not} vanish in general---the quartic pole produces curvature
(Computation~\ref{comp:virasoro-curvature}).
For affine Kac--Moody algebras ($\widehat{\fg}_k$ with double pole $\kappa_{ij}$
\emph{and} simple pole $f^{ab}{}_c J^c$), the simple pole contribution
$a_{(0)}b \neq 0$ gives nonzero residues, so the maximal-form
elements are not automatically cycles.
\end{remark}


\begin{proposition}[Kac--Moody acyclicity at generic level; \ClaimStatusProvedHere]
\label{prop:km-generic-acyclicity}
\index{Kac--Moody algebra!acyclicity at generic level}
For any simple Lie algebra $\fg$ and generic level $k \neq -h^\vee$,
the bar-twisted complex $\B(\hat{\fg}_k) \otimes_\tau V_k$ is acyclic:
\[
H_n(\B(\hat{\fg}_k) \otimes_\tau V_k) = \begin{cases}
\bC & n = 0, \\
0 & n > 0.
\end{cases}
\]
At the critical level $k = -h^\vee$, new homology classes appear,
corresponding to the Feigin--Frenkel center $\mathfrak{z}(\hat{\fg}_{-h^\vee})$.
\end{proposition}

\begin{proof}
Define the PBW filtration $F_p$ by the number of generators applied
to the vacuum.  The associated graded is
$\mathrm{gr}_F(\B(\hat{\fg}_k) \otimes_\tau V_k) \cong
\B(\fg[t^{-1}]) \otimes S(\fg[t^{-1}])$,
the bar complex of the loop algebra tensored with the symmetric algebra.
The resulting spectral sequence has
\[
E_1^{p,q} = H_p(\fg;\, S^q(\fg[t^{-1}]))
\]
(Lie algebra homology with polynomial coefficients).
By the Whitehead lemmas, $H_n(\fg; M) = 0$ for $n > 0$ and any
rational $\fg$-module~$M$, using the semisimplicity of~$\fg$.
The spectral sequence collapses at~$E_1$, giving acyclicity.

At $k = -h^\vee$, the vacuum module becomes reducible (the Sugawara
tensor is undefined), and the PBW filtration degenerates.
Computation~\ref{comp:km-acyclic} verifies the argument explicitly
for $\fg = \mathfrak{sl}_2$.
\end{proof}


\begin{proposition}[$\mathcal{W}_3$ vacuum leakage dichotomy; \ClaimStatusProvedHere]
\label{prop:w3-vacuum-dichotomy}
\index{vacuum leakage!$\mathcal{W}_3$ dichotomy}
\index{W3 algebra@$\mathcal{W}_3$!vacuum leakage}
In the degree-$2$ bar differential of the $\mathcal{W}_3$ algebra,
only the diagonal generator pairs produce vacuum leakage:
\begin{center}
\renewcommand{\arraystretch}{1.2}
\begin{tabular}{lcc}
\textbf{Pair} & \textbf{Vacuum leakage} & \textbf{Mechanism} \\
\hline
$T \otimes T$ & $c/2$ & quartic pole $T_{(3)}T = c/2$ \\
$W \otimes W$ & $c/3$ & sixth-order pole $W_{(5)}W = c/3$ \\
$T \otimes W$ & $0$ & pole order $\leq 2$; need order $\geq 5$ \\
$W \otimes T$ & $0$ & pole order $\leq 2$; need order $\geq 5$
\end{tabular}
\end{center}
Consequently, the curvature of the $\mathcal{W}_3$ bar complex decomposes
into two independent channels, both proportional to~$c$ and both vanishing
simultaneously if and only if $c = 0$.
The ratio $m_0^{(W)}/m_0^{(T)} = 2/3$ is independent of the level~$k$.
\end{proposition}

\begin{proof}
A vacuum contribution from $D(a \otimes b \otimes \eta_{12})
= \sum_{n \geq 0} a_{(n)}b$ requires $a_{(n_0)}b \in \bC \cdot |0\rangle$
for some~$n_0$.  The conformal weight condition forces
$n_0 = \mathrm{wt}(a) + \mathrm{wt}(b) - 1$.

\emph{Diagonal pairs.}
For $T \otimes T$: $n_0 = 2 + 2 - 1 = 3$, and $T_{(3)}T = c/2 \neq 0$.
For $W \otimes W$: $n_0 = 3 + 3 - 1 = 5$, and $W_{(5)}W = c/3 \neq 0$.

\emph{Off-diagonal pairs.}
For $T \otimes W$: $n_0 = 2 + 3 - 1 = 4$.  But $W$ is primary of
weight~$3$ under the Virasoro, so $T_{(n)}W = 0$ for $n \geq 2$
(the $TW$ OPE has pole order at most~$2$).  In particular $T_{(4)}W = 0$.
For $W \otimes T$: $n_0 = 4$, and $W_{(n)}T = 0$ for $n \geq 2$ by
skew-symmetry (Computation~\ref{comp:w3-nthproducts}).

The curvature ratio follows from the OPE normalizations:
$(c/3)/(c/2) = 2/3$, independent of~$k$.
Computation~\ref{comp:w3-bar-degree2} provides the complete verification.
\end{proof}

\begin{remark}[Generalization to $\mathcal{W}_N$]\label{rem:wn-vacuum-channels}
For $\mathcal{W}_N$ with generators $W^{(d_i)}$ of weights $d_1 = 2 < d_2 < \cdots < d_r$
(where $r = N-1$), the vacuum leakage dichotomy generalizes: the pair
$W^{(d_i)} \otimes W^{(d_j)}$ produces vacuum leakage if and only if
$i = j$ (i.e., the pair is diagonal).
This follows from the primary field condition: if $d_i \neq d_j$,
then the OPE pole order is at most $\max(d_i, d_j)$, while the vacuum
condition requires pole order $d_i + d_j - 1 > \max(d_i, d_j)$ when
$\min(d_i, d_j) \geq 2$.  Thus the curvature of any $\mathcal{W}_N$
bar complex decomposes into $r$ independent channels, one per generator.
\end{remark}


